% ============================================================
%  Electromagnetics — TikZ / PGFPlots styles for the figures
%  Roles: field blue = the field / the main object; copper = sources, charges, the second
%  object or the point of interest; teal = a third curve or a region; slate = guides.
%  Labels are mathematics in the math font (language neutral); words go through \ELL.
% ============================================================
\tikzset{
  ELax/.style={-{Stealth[length=5.5pt,width=4.2pt]},line width=0.6pt,ELink!80},
  ELgridl/.style={line width=0.3pt,ELgrid},
  ELcurve/.style={line width=1.15pt,ELcField},
  ELcurve2/.style={line width=1.15pt,ELcSource},
  ELcurve3/.style={line width=1.15pt,ELcThird},
  ELthin/.style={line width=0.6pt,ELcField!85},
  ELdash/.style={dashed,line width=0.55pt,ELcGuide},
  ELvec/.style={-{Stealth[length=6pt,width=4.6pt]},line width=1.05pt,ELcField},
  ELvec2/.style={-{Stealth[length=6pt,width=4.6pt]},line width=1.05pt,ELcSource},
  ELvec3/.style={-{Stealth[length=6pt,width=4.6pt]},line width=1.05pt,ELcThird},
  ELfieldline/.style={line width=0.75pt,ELcField,
    decoration={markings,mark=at position 0.5 with {\arrow{Stealth[length=5pt,width=4pt]}}},postaction=decorate},
  ELsurface/.style={fill=ELcField!10,draw=ELcField!70,line width=0.7pt},
  ELregion/.style={fill=ELcThird!12,draw=none},
  ELcharge/.style={circle,fill=ELcSource,inner sep=0pt,minimum size=5pt},
  ELchargeN/.style={circle,fill=ELcField,inner sep=0pt,minimum size=5pt},
  ELpt/.style={circle,fill=ELcSource,inner sep=0pt,minimum size=4.2pt},
  ELhole/.style={circle,fill=white,draw=ELcSource,line width=0.8pt,inner sep=0pt,minimum size=4.6pt},
  ELlab/.style={font=\small,text=ELink,inner sep=1.5pt},
  ELlabs/.style={font=\footnotesize,text=ELink,inner sep=1pt},
  ELlabg/.style={font=\footnotesize,text=ELink,inner sep=1.2pt,fill=white,fill opacity=0.85,text opacity=1},
  ELtag/.style={font=\small\itshape,text=ELcGuide,inner sep=2pt},
  ELmid/.style={decoration={markings,mark=at position 0.5 with {\arrow{Stealth[length=5pt]}}},postaction=decorate},
}
% a vector into the page (⊗) and out of the page (⊙) at a point
\newcommand{\ELinto}[2][ELcField]{\draw[#1,line width=0.7pt] #2 circle (3.2pt);
  \draw[#1,line width=0.7pt] ($#2+(-2.2pt,-2.2pt)$)--($#2+(2.2pt,2.2pt)$) ($#2+(-2.2pt,2.2pt)$)--($#2+(2.2pt,-2.2pt)$);}
\newcommand{\ELoutof}[2][ELcField]{\draw[#1,line width=0.7pt] #2 circle (3.2pt); \fill[#1] #2 circle (1.1pt);}
% rectangular axes: \ELaxes{xmin}{xmax}{ymin}{ymax}{x label}{y label}
\newcommand{\ELaxes}[6]{%
  \draw[ELax] (#1,0) -- (#2,0) node[ELlab,right]{#5};
  \draw[ELax] (0,#3) -- (0,#4) node[ELlab,above]{#6};}
% a right-handed 3-D frame at a point: \ELframe{origin}{length}{x}{y}{z}
\newcommand{\ELframe}[5]{%
  \draw[ELax] #1 -- ++(0,0,#2) node[ELlab,anchor=north east]{#3};
  \draw[ELax] #1 -- ++(#2,0,0) node[ELlab,anchor=west]{#4};
  \draw[ELax] #1 -- ++(0,#2,0) node[ELlab,anchor=south]{#5};}
\pgfplotsset{ELplot/.style={
  axis lines=middle,axis line style={ELink!80,-{Stealth[length=5pt]}},
  tick style={ELink!60},tick label style={font=\footnotesize,text=ELink,fill=white,fill opacity=0.85,
    text opacity=1,inner sep=1.3pt},
  label style={font=\small,text=ELink},
  every axis x label/.style={at={(ticklabel* cs:1.0)},anchor=west,xshift=3pt},
  every axis y label/.style={at={(ticklabel* cs:1.0)},anchor=south,yshift=3pt},
  legend style={font=\footnotesize,draw=ELhair,fill=white,fill opacity=0.9,text opacity=1},
  cycle list={{ELcField},{ELcSource},{ELcThird},{ELcGuide,dashed}},
  grid=none,width=9cm,height=5.2cm,every axis plot/.append style={line width=1.15pt}}}
% the oblique projection of the 3-D drawings (x toward the reader, lower left; y right; z up)
\tikzset{ELoblique/.style={x={(-0.5cm,-0.42cm)},y={(1cm,0cm)},z={(0cm,1cm)}}}
% a concept tree (the outline slides): boxes and the joining lines
\tikzset{
  ELtreebox/.style={draw=ELcField!45,fill=ELblueT,rounded corners=2pt,line width=0.6pt,align=center,
    font=\small\bfseries,text=ELink,inner sep=4pt,minimum height=0.95cm,
    execute at begin node={\hyphenpenalty=10000\exhyphenpenalty=10000}},
  ELtreeboxR/.style={ELtreebox,draw=ELcSource!45,fill=ELredT},
  ELtreeroot/.style={ELtreebox,draw=ELslate,fill=ELslateT},
  ELtreeon/.style={fill=ELcField,draw=ELcField,text=white},
  ELtreeonR/.style={fill=ELcSource,draw=ELcSource,text=white},
  ELtreeline/.style={line width=0.6pt,draw=ELcField!55},
  ELtreelineR/.style={line width=0.6pt,draw=ELcSource!55},
  ELtreelineG/.style={line width=0.6pt,draw=ELslate!70}}
% the cyclic order of three unit vectors: \ELcycle{top}{right}{left} (clockwise arrows)
\newcommand{\ELcycle}[3]{%
  \begin{tikzpicture}[scale=0.8]
    \node[ELlab] (t) at (90:1) {#1}; \node[ELlab] (r) at (-30:1) {#2}; \node[ELlab] (l) at (210:1) {#3};
    \draw[ELcField,line width=0.8pt,-{Stealth[length=5pt]}] (72:1) arc (72:-12:1);
    \draw[ELcField,line width=0.8pt,-{Stealth[length=5pt]}] (-48:1) arc (-48:-132:1);
    \draw[ELcField,line width=0.8pt,-{Stealth[length=5pt]}] (192:1) arc (192:112:1);
  \end{tikzpicture}}
% a resistor (or a reluctance) drawn upward from the point where it is placed, 1.15 units long
\tikzset{ELresV/.pic={\draw[ELink,line width=0.8pt,line join=round] (0,0)--(0,0.15)--(0.12,0.22)--(-0.12,0.36)--(0.12,0.5)
  --(-0.12,0.64)--(0.12,0.78)--(-0.12,0.92)--(0,0.99)--(0,1.15);}}
